\subsection*{一步一步建立直观}

\paragraph{1. 把复杂曲线拆成可单独追踪的成分}

先在初始时刻观察两条正弦曲线及其和。局部看起来复杂的起伏，可能只是少数简单模式叠加。傅里叶展开把这种想法推广到一组频率，并用投影计算系数。右图只取两个分量，目的是让每个成分的贡献可以被直接辨认。



\paragraph{2. 同一种分解可以服务不同问题}

推进时间后，每个系数按热方程的规则衰减，基函数本身没有更换。若研究声音，则常把系数看成不同频率成分；若研究旋转向量，则用复指数表示同样的频率结构。本窗口展示热扩散用途，并不把时间衰减当作所有傅里叶问题的固有属性。



\paragraph{3. 增加频率并不等于增加任意细节}

改变第二个模式的频率，比较它与总曲线的关系。一个固定频率、振幅和相位只能添加特定的周期形状；重建一般目标函数需要很多系数，并受收敛方式影响。跳跃附近的行为尤其要单独分析，不能把有限平滑曲线的视觉效果直接推广到所有函数。
